%% kapa: exam=1 round=31 year=2020 period=1 form=A subject=civil_law subject_ko=민법 seq=4 number=4
%% source: Q-Net 감정평가사 기출문제 (cst00309) artlSeq=5207219 file=extracted/1008/5207219/A형/A형/2020년도 제31회 감정평가사 1차 시험_1교시_A형.pdf
%% conversion: kordoc 4.24.51 -> tools/gen1.py (auto); review=proofread-sonnet; flags=none
\begin{question}{4}
부재자의 재산관리에 관한 설명으로 옳지 않은 것은? (다툼이 있으면 판례에 따름)
\begin{choices}
  \item 부재자가 스스로 위임한 재산관리인에게 재산처분권까지 준 경우에도 그 재산관리인은 재산처분에 법원의 허가를 얻어야 한다.
  \item 재산관리인의 권한초과행위에 대한 법원의 허가결정은 기왕의 처분행위를 추인하는 방법으로도 할 수 있다.
  \item 재산관리인이 소송절차를 진행하던 중 부재자에 대한 실종선고가 확정되면 그 재산관리인의 지위도 종료한다.
  \item 생사불명의 부재자를 위하여 법원이 선임한 재산관리인은 그가 부재자의 사망을 확인한 때에도 선임결정이 취소되지 않으면 계속 권한을 행사할 수 있다.
  \item 생사불명의 부재자에 대하여 실종이 선고되더라도 법원이 선임한 재산관리인의 처분행위에 근거한 등기는 그 선임결정이 취소되지 않으면 적법하게 마친 것으로 추정된다.
\end{choices}
\end{question}
\answer{1}
